%% ============================================================
%%  Chapter 22 --- Durable Execution & Long-Running Workflows
%%  Production-Grade Agentic AI: LangGraph + Python Frameworks
%%  Dependencies: langgraph >= 0.2.0,
%%                langgraph-checkpoint-postgres >= 0.1.0,
%%                langchain-core >= 0.3.0,
%%                pydantic >= 2.7.0,
%%                fastapi >= 0.115.0,
%%                asyncpg >= 0.29.0,
%%                celery[redis] >= 5.4.0,
%%                langgraph-sdk >= 0.1.0
%% ============================================================

\chapterquote{A system that loses its work when the process dies is not
an agent. It is a script with ambitions.}{anon}

The production failure this chapter prevents is silent work loss. The
approval workflow of Chapter~21 records every decision in an
\code{InMemorySaver} checkpointer---an in-process \code{defaultdict}
that exists only as long as the Python interpreter lives. A \code{SIGTERM}
signal from a Kubernetes node eviction, an unhandled exception in an
adjacent coroutine, or a routine deployment rollout terminates the process
and erases every checkpoint. When the process restarts, \code{ainvoke}
on the same \code{thread\_id} finds no prior state, silently begins a
fresh run, and overwrites the execution history without a warning. An
agent workflow that required ten minutes, three LLM calls, and two human
approvals to reach its third milestone must be restarted from scratch. At
scale---hundreds of concurrent long-running workflows, each spanning
minutes to hours---this loss is not an edge case. It is a guarantee that
every infrastructure event resets every active workflow simultaneously.

The solution has four layers. The first is a \emph{persistence layer}:
\code{AsyncPostgresSaver} replaces \code{InMemorySaver} as the graph's
checkpointer, writing each state transition to a Postgres table that
survives process termination. The second is a \emph{recovery layer}: a
\code{RecoveryNode} placed at the graph's entry point queries the
checkpointer at startup, locates the highest committed milestone for the
current thread, and routes execution forward from that point rather than
from \code{START}. The third is a \emph{dispatch layer}: a Celery task
queue backed by Redis decouples the HTTP request lifecycle from graph
execution, ensuring that a graph run deferred to a worker continues
independently of the originating request and is retried automatically if
the worker crashes before completion. The fourth is a \emph{managed
alternative}: LangGraph Platform provides the first three layers as
infrastructure, requiring only a single SDK call in deployments where
self-hosting is not a requirement.

The sections are ordered by the dependency each layer has on the one
before it.
Section~\ref{sec:durable-postgres} establishes the persistence layer
first because the recovery layer has no durable checkpoints to query
without it.
Section~\ref{sec:durable-recovery} adds the \code{RecoveryNode} after
persistence is established, because the node's contract---find the
highest completed milestone and resume from its
\code{CheckpointTuple.config}---requires
\code{AsyncPostgresSaver.alist()} to be available.
Section~\ref{sec:durable-celery} introduces the dispatch layer after
recovery, because a Celery task that invokes the graph inherits both
the persistence and recovery guarantees; the task's own retry contract
is complementary to, not a replacement for, checkpoint-based recovery.
Section~\ref{sec:durable-platform} presents the managed alternative
after the self-hosted stack is fully established, so the trade-offs
between the two models are concrete.
Section~\ref{sec:durable-full} provides the integration test that
exercises all four layers together by spawning three sequential
subprocesses on the same \code{thread\_id} and asserting that execution
resumes from the correct milestone after each termination.

This chapter introduces seven terms defined in
Table~\ref{tab:ch22-terms}.

\begin{table}[h]
\centering
\caption{Domain terms introduced in Chapter~22}
\label{tab:ch22-terms}
\begin{tabularx}{\textwidth}{lX}
\toprule
\textbf{Term} & \textbf{Definition} \\
\midrule
\code{AsyncPostgresSaver} &
  \code{from langgraph.checkpoint.postgres.aio import
  AsyncPostgresSaver}; a \code{BaseCheckpointSaver} implementation
  that persists checkpoints to a Postgres \code{checkpoints} table;
  constructed via \code{async with
  AsyncPostgresSaver.from\_conn\_string(dsn) as saver}; schema
  initialised by \code{await saver.setup()} which runs idempotent
  \code{CREATE TABLE IF NOT EXISTS} DDL on every call. \\
\code{InMemorySaver} &
  \code{from langgraph.checkpoint.memory import InMemorySaver};
  the canonical name (as of LangGraph~0.2) for the in-process
  checkpointer used in Chapters~19--21 under the alias
  \code{MemorySaver}; all checkpoints are lost on process exit;
  never suitable for production long-running workflows. \\
\code{alist()} &
  \code{checkpointer.alist(config,~*,~filter,~before,~limit)};
  returns \code{AsyncIterator[CheckpointTuple]}; yields checkpoints
  newest-first; each \code{CheckpointTuple} has \code{.checkpoint}
  (a dict whose \code{"channel\_values"} key holds the full state),
  \code{.metadata} (\code{source}, \code{step}, \code{writes}),
  \code{.config} (with \code{checkpoint\_id}), and
  \code{.parent\_config}. \\
\code{CheckpointTuple} &
  \code{from langgraph.checkpoint.base import CheckpointTuple};
  the unit yielded by \code{alist()}; has no \code{.values} or
  \code{.next} attributes---those belong to \code{StateSnapshot};
  \code{.checkpoint["channel\_values"]} holds the state dict;
  \code{.metadata["writes"]} is \code{dict[node\_name,
  node\_return\_value]}, keyed by the registered node name that
  produced the checkpoint. \\
\code{RecoveryNode} &
  A graph node placed immediately after \code{START} with a
  conditional edge; queries \code{alist()} for the highest committed
  milestone checkpoint; routes to the appropriate resume node when a
  milestone is found, or to \code{research\_complete} when no
  milestone exists; performs no LLM calls or external writes. \\
\code{MILESTONE\_NODES} &
  A \code{frozenset[str]} of \emph{registered node names} that
  represent safe resume points; a checkpoint is a milestone checkpoint
  when \code{metadata["writes"]} contains at least one key from
  \code{MILESTONE\_NODES}; node names in this set must exactly match
  the strings passed to \code{add\_node()} in \code{build\_durable\_graph()}. \\
\code{acks\_late} &
  A Celery task configuration flag; when \code{True}, the task
  message is acknowledged (removed from the queue) only after the
  task function returns, not when the worker receives it; prevents
  silent task loss when a worker process is killed mid-execution;
  must be set in both the \code{@task} decorator and
  \code{celery\_app.conf.update()} to take effect in all broker
  configurations. \\
\bottomrule
\end{tabularx}
\end{table}

\medskip
\noindent\textbf{Source layout for this chapter:}
\begin{verbatim}
chapter22/
  checkpointer.py        # §22.1 — lifespan_checkpointer,
                         #          get_checkpointer
  state.py               # §22.5 — DurableWorkflowState
  recovery.py            # §22.2 — MILESTONE_NODES,
                         #          find_resume_checkpoint,
                         #          recovery_node,
                         #          route_after_recovery
  graph.py               # §22.5 — build_durable_graph,
                         #          node functions
  tasks.py               # §22.3 — celery_app,
                         #          run_workflow_task
  platform_runner.py     # §22.4 — create_background_run,
                         #          poll_run_status
  worker.py              # §22.5 — subprocess entry point,
                         #          asyncio SIGTERM handler
  test_three_restarts.py # §22.5 — three-subprocess harness
\end{verbatim}


%% ------------------------------------------------------------
\section{Durable State Outlives the Process That Wrote It}
\label{sec:durable-postgres}
%% ------------------------------------------------------------

The design problem is the gap between when a checkpoint is written and
when the process that wrote it is expected to still be running.
\code{InMemorySaver}---the checkpointer used in Chapters~19 through~21
for its zero-configuration convenience---stores every checkpoint in a
Python \code{defaultdict} that exists only in the heap of the current
process. The moment that process exits, for any reason, every checkpoint
disappears. The naive mitigation is to catch \code{SIGTERM} and flush
state before exiting; this works for graceful shutdowns but fails for
\code{SIGKILL}, out-of-memory kills, and hardware failures, which are
exactly the events a production container scheduler produces most
frequently. The governing principle is that a checkpoint must be durable
the instant it is written---the process that wrote it must be able to
die the next millisecond without data loss.

\code{AsyncPostgresSaver} achieves this by writing each checkpoint as a
row in a Postgres \code{checkpoints} table inside the same transaction
that the Pregel scheduler uses to advance the graph. Postgres's
write-ahead log ensures that a committed row survives process termination,
node eviction, and instance replacement without additional application-level
flushing. The checkpointer is constructed via
\code{async with AsyncPostgresSaver.from\_conn\_string(dsn) as saver}:
\code{from\_conn\_string} is an \emph{async context manager}, meaning the
\code{saver} object is only valid inside the \code{async with} block; the
block must remain open for the lifetime of the process, not just for the
duration of a single graph invocation. The correct placement is inside a
FastAPI \code{lifespan} context, which opens on worker startup and closes
on graceful shutdown.

\code{setup()} must be called once per process startup, not once per
graph invocation. It executes idempotent \code{CREATE TABLE IF NOT EXISTS}
DDL that creates the \code{checkpoints}, \code{checkpoint\_blobs}, and
\code{checkpoint\_migrations} tables when they do not yet exist. Calling
it on every invocation is harmless but wasteful; calling it never means
the first \code{ainvoke} on a freshly deployed instance raises a
\code{UndefinedTableError} before writing the first checkpoint.

\begin{lstlisting}[language=Python,
  caption={AsyncPostgresSaver as a lifespan-scoped context manager;
           get\_checkpointer factory with startup guard},
  label={lst:durable-checkpointer}]
# chapter22/checkpointer.py
from contextlib import asynccontextmanager
from typing import AsyncIterator
from langgraph.checkpoint.postgres.aio import AsyncPostgresSaver

_saver: AsyncPostgresSaver | None = None

@asynccontextmanager
async def lifespan_checkpointer(
    dsn: str,
) -> AsyncIterator[AsyncPostgresSaver]:
    global _saver
    async with AsyncPostgresSaver.from_conn_string(dsn) as saver:
        await saver.setup()
        _saver = saver
        yield saver
        _saver = None

def get_checkpointer() -> AsyncPostgresSaver:
    if _saver is None:
        raise RuntimeError(
            "Checkpointer not initialised. "
            "Use lifespan_checkpointer() in FastAPI lifespan."
        )
    return _saver
\end{lstlisting}

Two decisions in Listing~\ref{lst:durable-checkpointer} deserve
explanation.
\code{\_saver} is a module-level variable rather than a
dependency-injection singleton because \code{AsyncPostgresSaver} holds
an open connection pool that must be shared across all concurrent
requests; constructing a new instance per request would exhaust the
Postgres connection limit under load and would re-parse the DSN string---
which may contain credentials---on every invocation. The
\code{get\_checkpointer()} guard raises \code{RuntimeError} rather than
silently returning \code{InMemorySaver()} as a fallback: a fallback would
make the production failure mode (silent work loss) undetectable in
staging environments that happen to start the checkpointer correctly; an
explicit error surfaces the misconfiguration at the first graph invocation
rather than at the first process restart when all work is already lost.

\begin{notebox}
Chapter~3~\S3.2 established \code{settings.POSTGRES\_DSN} as the
canonical DSN source for this book's infrastructure. Pass it directly to
\code{lifespan\_checkpointer(settings.POSTGRES\_DSN)}. Do not open a
second \code{asyncpg} connection pool independently of the one managed by
\code{AsyncPostgresSaver}; the checkpointer manages its own internal pool
from the DSN string, sized separately from the application's query pool.
\end{notebox}

The wiring into FastAPI is a single \code{lifespan} context that nests
\code{lifespan\_checkpointer} before yielding control to the application:

\begin{lstlisting}[language=Python,
  caption={FastAPI lifespan wiring: checkpointer opens on startup and
           closes on shutdown},
  label={lst:durable-lifespan}]
# chapter22/app.py  (excerpt)
from contextlib import asynccontextmanager
from fastapi import FastAPI
from chapter22.checkpointer import lifespan_checkpointer
from chapter22.settings import settings   # Ch.14 §14.3 AppSettings

@asynccontextmanager
async def lifespan(app: FastAPI):
    async with lifespan_checkpointer(settings.POSTGRES_DSN):
        yield   # application serves requests with checkpointer open

app = FastAPI(lifespan=lifespan)
\end{lstlisting}

Two decisions in Listing~\ref{lst:durable-lifespan} deserve explanation.
\code{lifespan\_checkpointer} is composed with the \code{app}'s lifespan
via nesting rather than through a shared lifespan utility: FastAPI accepts
exactly one \code{lifespan} callable per application instance, so all
startup and shutdown logic must chain through that single callable;
nesting \code{async with} blocks is the correct composition pattern and
avoids the signal-handling conflicts that arise when multiple lifespan
managers independently register \code{SIGTERM} handlers. The \code{yield}
is bare---it carries no application state---because \code{get\_checkpointer()}
provides the module-level access point; any handler that needs the
checkpointer calls \code{get\_checkpointer()} directly, which avoids
requiring every route handler to declare an injection dependency for a
resource that most handlers never touch.

The persistence layer is now in place: every checkpoint written by a
graph run is durable from the moment of the \code{await}. The remaining
problem is what happens when a new process starts and calls \code{ainvoke}
on a \code{thread\_id} that already has prior durable state---the graph
must resume from the highest completed milestone rather than restarting
from \code{START}.


%% ------------------------------------------------------------
\section{A Restarted Process Must Resume, Not Restart}
\label{sec:durable-recovery}
%% ------------------------------------------------------------

The design problem is the default behaviour of \code{ainvoke} when called
on a thread that already has checkpoints. The Pregel scheduler loads the
latest checkpoint for the given \code{thread\_id} and resumes from
\code{StateSnapshot.next}---the set of nodes that had not yet completed
when execution was last interrupted. If the process was killed cleanly
between node completions, \code{.next} points to the correct resume node
and no special handling is required. But if the process was killed
mid-node or mid-checkpoint-write, the checkpoint is either absent or
records a node that has partially executed its side effects. The naive
fix---catching \code{SIGTERM} and hoping the checkpoint was already
flushed---is a race condition. The governing principle is that a graph
must resume only from a \emph{known safe milestone}: a node whose
completion checkpoint is guaranteed consistent, not whatever the
checkpoint table happens to record as the latest \code{.next}.

\emph{Milestone nodes} are nodes whose completion represents a
semantically coherent, externally observable unit of work: a research
phase completed, a draft written, a review passed. Non-milestone nodes---
routing helpers, state-mutation steps, intermediate sub-nodes---may
produce incomplete checkpoints when a process dies mid-execution; re-running
them from the top is safe because they are idempotent by the contract from
Chapter~21~\S21.1. The \code{RecoveryNode} is placed immediately after
\code{START}: it queries \code{AsyncPostgresSaver.alist()} for the highest
committed milestone checkpoint, then routes execution directly to the node
that follows that milestone, bypassing all nodes that have already
completed safely.

\code{alist(config)} returns an \code{AsyncIterator[CheckpointTuple]}
yielding newest-first. Each \code{CheckpointTuple.metadata["writes"]} is
a \code{dict[str,~Any]} mapping the \emph{registered node name} that
produced the checkpoint to that node's return value. Checking whether any
key in \code{metadata["writes"]} belongs to \code{MILESTONE\_NODES}
identifies milestone checkpoints as a pure metadata operation---no state
deserialisation is required.

\begin{warningbox}
Do not use \code{graph.get\_state\_history(config)} inside
\code{RecoveryNode}. \code{get\_state\_history} returns
\code{Iterator[StateSnapshot]}, which deserialises every checkpoint's
full state values as it iterates. \code{alist()} reads only metadata
until the caller accesses \code{.checkpoint["channel\_values"]}; in a
thread with hundreds of checkpoints this is the difference between a
recovery that takes milliseconds and one that deserialises megabytes of
state before finding the first milestone.
\end{warningbox}

One constraint governs the relationship between \code{MILESTONE\_NODES}
and the graph's node registry: every string in \code{MILESTONE\_NODES}
must \emph{exactly} match the name passed to \code{b.add\_node(name,~fn)}
in \code{build\_durable\_graph()}. \code{metadata["writes"]} is keyed by
the registered node name, not by the function name. A mismatch between
the frozenset and the registry produces a recovery scan that never finds a
milestone and always restarts from \code{START}---silently and without
any error.

\begin{lstlisting}[language=Python,
  caption={MILESTONE\_NODES frozenset and alist() newest-first
           milestone scan returning CheckpointTuple},
  label={lst:durable-find-checkpoint}]
# chapter22/recovery.py
from langgraph.checkpoint.postgres.aio import AsyncPostgresSaver
from langgraph.checkpoint.base import CheckpointTuple
from langchain_core.runnables import RunnableConfig

# Names must exactly match b.add_node() calls in build_durable_graph().
MILESTONE_NODES: frozenset[str] = frozenset({
    "research_complete",
    "draft_complete",
    "review_complete",
})

async def find_resume_checkpoint(
    checkpointer: AsyncPostgresSaver,
    thread_id:    str,
) -> CheckpointTuple | None:
    config: RunnableConfig = {
        "configurable": {"thread_id": thread_id}
    }
    # Yields newest-first; first match is the highest milestone.
    async for tup in checkpointer.alist(config):
        writes: dict = tup.metadata.get("writes") or {}
        if MILESTONE_NODES.intersection(writes.keys()):
            return tup   # closes iterator immediately
    return None
\end{lstlisting}

Two decisions in Listing~\ref{lst:durable-find-checkpoint} deserve
explanation.
The function returns \code{CheckpointTuple~|~None} rather than
\code{StateSnapshot}: \code{CheckpointTuple.config} carries the
\code{checkpoint\_id} required to resume from that exact checkpoint,
whereas obtaining the equivalent from a \code{StateSnapshot} would
require first calling \code{graph.aget\_state(tup.config)}, which
deserialises the full state; returning the raw tuple defers
deserialisation to the caller and keeps the type unambiguous.
\code{return tup} inside the \code{async for} loop closes the
\code{alist()} iterator immediately: \code{AsyncPostgresSaver.alist()}
uses server-side cursors that hold a Postgres transaction open until
the iterator is exhausted or closed; an unreturned iterator that is
simply abandoned holds that transaction open until the garbage collector
finalises it, which under high concurrency exhausts the connection pool.

With the milestone lookup established, the \code{RecoveryNode} is a
thin wrapper that calls \code{find\_resume\_checkpoint()} and writes a
\code{stage} routing signal into state. The checkpointer is passed as a
parameter to \code{build\_durable\_graph()} and threaded into
\code{recovery\_node} via a closure, avoiding any dependency on the
module-level \code{get\_checkpointer()} singleton inside node bodies---
node bodies should be free functions that receive all dependencies from
the outside.

\begin{lstlisting}[language=Python,
  caption={RecoveryNode factory, stage extraction from milestone name,
           and routing function for the conditional edge},
  label={lst:durable-recovery-node}]
# chapter22/recovery.py  (continued)
from chapter22.state import DurableWorkflowState

def make_recovery_node(checkpointer: AsyncPostgresSaver):
    async def recovery_node(
        state:  DurableWorkflowState,
        config: RunnableConfig,
    ) -> dict:
        thread_id: str = (
            (config.get("configurable") or {}).get("thread_id", "")
        )
        tup = await find_resume_checkpoint(checkpointer, thread_id)
        if tup is None:
            return {"stage": "init"}   # no prior milestone: start fresh
        milestone = next(
            k for k in (tup.metadata.get("writes") or {})
            if k in MILESTONE_NODES
        )
        # "research_complete" -> "research", etc.
        stage = milestone.replace("_complete", "")
        return {"stage": stage}
    return recovery_node

def route_after_recovery(state: DurableWorkflowState) -> str:
    return {
        "init":     "research_complete",
        "research": "draft_complete",
        "draft":    "review_complete",
        "review":   "synthesise_node",
    }.get(state["stage"], "research_complete")
\end{lstlisting}

Three decisions in Listing~\ref{lst:durable-recovery-node} deserve
explanation.
\code{make\_recovery\_node} is a factory that returns a closure over
\code{checkpointer} rather than calling \code{get\_checkpointer()}
directly inside the node body: this makes the checkpointer dependency
explicit at graph-construction time, keeps \code{recovery\_node} a pure
function of its arguments, and allows tests to inject a mock checkpointer
without patching a module global. \code{recovery\_node} writes only
\code{\{"stage":~stage\}} rather than replaying the full state from
\code{tup.checkpoint["channel\_values"]}: the Pregel scheduler loads the
correct full state from the checkpoint identified by
\code{tup.config["configurable"]["checkpoint\_id"]} when it re-enters the
resume node; writing the full state here would invoke all reducers
unconditionally and risk overwriting values written by nodes that completed
after the milestone. \code{route\_after\_recovery} is a module-level free
function rather than a method because LangGraph's
\code{add\_conditional\_edges} accepts a callable; a free function with a
clear \code{DurableWorkflowState} signature is independently testable
without any graph infrastructure.

The recovery layer ensures a restarted process resumes from the correct
milestone. The remaining problem is upstream: what guarantees that the
original \code{ainvoke} call retries at all when the process that received
the HTTP request is killed before the graph run completes?


%% ------------------------------------------------------------
\section{Long-Running Work Must Outlive the Request That Triggered It}
\label{sec:durable-celery}
%% ------------------------------------------------------------

The design problem is the coupling between the HTTP request lifecycle and
graph execution duration. A graph run spanning minutes is initiated by an
HTTP \code{POST} request. A synchronous response requires the client to
hold an open connection for the run's duration, which is impractical for
any run longer than a load balancer's idle timeout. The naive fix---using
\code{asyncio.create\_task()} to run the graph in the background and
responding immediately with 202---decouples the client from the response
but not the work from the process: if the worker process is killed while
the graph runs in a background coroutine, the coroutine is silently dropped
with no retry and no record in any queue. The governing principle is that a
graph run which must survive the process that triggered it requires a
\emph{durable task queue}: a system that persists the task message to an
external store before any worker picks it up, so that a worker crash
triggers automatic re-delivery to a healthy worker.

\emph{Celery with a Redis broker} is the standard Python implementation
of a durable task queue. When a caller invokes
\code{run\_workflow\_task.apply\_async()}, Celery serialises the task
arguments to JSON and writes the message to Redis before returning. A
worker process picks up the message and begins executing the task function.
The two durability guarantees are complementary and non-redundant: Celery
guarantees the task function runs at least once even if the first worker
dies mid-execution; \code{AsyncPostgresSaver} guarantees that if the
Celery retry re-enters the graph on the same \code{thread\_id}, the
\code{RecoveryNode} resumes from the last committed milestone rather than
restarting from \code{START}. Without Celery, the task may never retry.
Without the checkpointer, the retry restarts from \code{START}.

Celery workers use a synchronous \code{prefork} process pool by default.
The task function is an ordinary Python function, not a coroutine. To
call the async \code{ainvoke} method, the task function must open a new
event loop via \code{asyncio.run()}: this is the officially supported
bridge between Celery's synchronous execution model and the async
LangGraph runtime.

\begin{warningbox}
Do not use \code{celery[gevent]} with \code{asyncpg}. The \code{gevent}
pool monkey-patches Python's socket layer at the process level;
\code{asyncpg} uses \code{asyncio}'s native socket primitives and is not
compatible with gevent's patched equivalents. The combination produces
intermittent deadlocks under load that are difficult to reproduce in
testing. Use the default \code{prefork} pool with \code{asyncio.run()}.
\end{warningbox}

\begin{lstlisting}[language=Python,
  caption={Celery app instantiation, run\_workflow\_task with
           asyncio.run() bridge, self.retry() on transient failure},
  label={lst:durable-celery-task}]
# chapter22/tasks.py
import asyncio
from celery import Celery
from celery.utils.log import get_task_logger
from chapter22.settings import settings   # Ch.14 §14.3 AppSettings

celery_app = Celery(
    "chapter22",
    broker=settings.REDIS_URL,
    backend=settings.REDIS_URL,
)
logger = get_task_logger(__name__)

@celery_app.task(
    bind=True,
    name="chapter22.tasks.run_workflow",
    max_retries=3,
    default_retry_delay=60,
    acks_late=True,
)
def run_workflow_task(
    self,
    thread_id: str,
    task:      str,
) -> dict:
    async def _run() -> dict:
        from chapter22.checkpointer import lifespan_checkpointer
        from chapter22.graph import build_durable_graph
        from chapter22.state import DurableWorkflowState
        async with lifespan_checkpointer(settings.POSTGRES_DSN) as cp:
            graph = build_durable_graph(cp)
            cfg   = {"configurable": {"thread_id": thread_id}}
            # RecoveryNode routes: fresh state only used on init.
            return await graph.ainvoke(
                DurableWorkflowState(task=task, stage="init",
                    research_done=False, draft_done=False,
                    review_done=False, final_output=""),
                config=cfg,
            )
    try:
        return asyncio.run(_run())
    except Exception as exc:
        logger.warning("run_workflow_task transient failure: %s", exc)
        raise self.retry(exc=exc)
\end{lstlisting}

Three decisions in Listing~\ref{lst:durable-celery-task} deserve
explanation.
\code{lifespan\_checkpointer} is opened inside the async \code{\_run()}
helper rather than at module scope: each Celery \code{prefork} worker is
a forked child process; an \code{asyncpg} connection pool opened in the
parent process before \code{fork()} cannot be used in the child because
the pool's file descriptors and event loop reference are copied but the
underlying socket state is not safely shareable; opening the pool inside
\code{asyncio.run()}---which creates a fresh event loop per task
invocation---avoids this. Passing a fresh \code{DurableWorkflowState}
as input on every call is safe because \code{RecoveryNode} runs first:
when a prior checkpoint exists, \code{recovery\_node} writes the
appropriate \code{stage} and the Pregel scheduler resumes from the
milestone's checkpoint state, ignoring the initial-state values for all
fields that have been committed; when no prior checkpoint exists,
the initial state is used as-is. \code{self.retry(exc=exc)} re-raises
the exception as a Celery retry signal rather than returning a default
value: a task that catches an exception and returns normally signals
success to Celery and stops all further retries; wrapping the exception
in \code{self.retry()} preserves Celery's exponential backoff and
\code{max\_retries} enforcement.

\begin{lstlisting}[language=Python,
  caption={Celery configuration block: acks\_late enforcement,
           prefetch multiplier, result expiry},
  label={lst:durable-celery-conf}]
# chapter22/tasks.py  (continued)
celery_app.conf.update(
    task_serializer="json",
    result_serializer="json",
    accept_content=["json"],
    # Must also appear here; decorator acks_late alone is
    # insufficient in some broker/transport configurations.
    task_acks_late=True,
    # Prevents one worker reserving N long tasks and starving others.
    worker_prefetch_multiplier=1,
    # Bound Redis memory: expire stored results after one hour.
    result_expires=3600,
)
\end{lstlisting}

Two decisions in Listing~\ref{lst:durable-celery-conf} deserve
explanation.
\code{worker\_prefetch\_multiplier=1} is the most consequential setting
for long-running task queues: the default value of~4 allows each worker
to pre-fetch and reserve four tasks simultaneously; a worker executing a
ten-minute graph task while holding three reserved tasks prevents those
three tasks from being picked up by any other worker for the full ten
minutes, creating queue latency that grows linearly with task duration.
\code{result\_expires=3600} bounds Redis memory use: without expiry,
the result backend accumulates one stored result entry per completed task
indefinitely; at hundreds of tasks per day this grows to gigabytes within
weeks and is never automatically reclaimed by Redis's eviction policies
unless a global \code{maxmemory-policy} is configured separately.

The dispatch layer ensures graph runs are durable across worker crashes
and are retried automatically. For teams operating on LangGraph Platform,
these three layers---persistence, recovery, and dispatch---are provided
as managed infrastructure, reducing the full stack to a single API call.


%% ------------------------------------------------------------
\section{Managed Infrastructure Compresses Three Layers to One Call}
\label{sec:durable-platform}
%% ------------------------------------------------------------

The design problem for teams using LangGraph Platform is the operational
overhead of maintaining three independent infrastructure components: a
Postgres instance for checkpoints, a Redis instance for the Celery broker
and result backend, and a worker fleet with correct \code{acks\_late}
configuration. Each component requires capacity planning, connection pool
tuning, schema migrations, and monitoring. The Platform provisions all
three as managed services and exposes them through a single API primitive:
the \emph{background run}. The governing principle is that the infrastructure
concerns solved in Sections~\ref{sec:durable-postgres}
through~\ref{sec:durable-celery} do not disappear on Platform---they are
transferred to the operator---and understanding the self-hosted implementation
makes the Platform's behaviour and limitations comprehensible rather than
opaque.

\code{AsyncLangGraphClient} is the Python SDK entry point for the LangGraph
Platform API. A background run is created by calling
\code{client.runs.create()} with the \code{thread\_id}, \code{assistant\_id},
and input dict; the Platform queues the run on its internal worker fleet,
persists the task durably, and returns a \code{run} object immediately.
The \code{multitask\_strategy} parameter controls behaviour when a second
run is submitted for a thread that is already active: \code{"enqueue"}
queues it behind the current run; \code{"interrupt"} cancels the current
run and starts the new one; \code{"rollback"} cancels and reverts state
to the pre-run checkpoint; \code{"reject"} returns an error without
starting a new run. For long-running sequential workflows where each
invocation represents a distinct phase, \code{"enqueue"} is correct.

\begin{lstlisting}[language=Python,
  caption={LangGraph Platform background run creation and status
           polling via AsyncLangGraphClient},
  label={lst:durable-platform}]
# chapter22/platform_runner.py
from langgraph_sdk import AsyncLangGraphClient
from chapter22.settings import settings

async def create_background_run(
    client:       AsyncLangGraphClient,
    assistant_id: str,
    thread_id:    str,
    task:         str,
) -> str:   # returns run_id
    run = await client.runs.create(
        thread_id=thread_id,
        assistant_id=assistant_id,
        input={"task": task},
        multitask_strategy="enqueue",
    )
    return run["run_id"]

async def poll_run_status(
    client:    AsyncLangGraphClient,
    thread_id: str,
    run_id:    str,
) -> str:   # "pending"|"running"|"success"|"error"|"timeout"|"interrupted"
    run = await client.runs.get(thread_id, run_id)
    return run["status"]

def make_client() -> AsyncLangGraphClient:
    return AsyncLangGraphClient(
        url=settings.LANGGRAPH_API_URL,
        api_key=settings.LANGGRAPH_API_KEY,
    )
\end{lstlisting}

Two decisions in Listing~\ref{lst:durable-platform} deserve explanation.
\code{poll\_run\_status} returns the raw status string rather than mapping
it to a Celery-compatible enum: the caller---typically a FastAPI status
endpoint---presents the string to the client unchanged; a mapping layer
would diverge whenever the Platform adds a new status value
(e.g.,~\code{"interrupted"} was added after the initial release), requiring
a code change to handle a state that the caller would otherwise propagate
correctly by default. \code{make\_client()} is a factory function rather
than a module-level singleton: \code{AsyncLangGraphClient} holds an
\code{httpx.AsyncClient} that must be closed after use; a factory makes the
client lifetime explicit at the call site and avoids the shared-mutable-state
problem that a module singleton introduces in test environments where multiple
clients with different credentials must coexist.

\begin{notebox}
Platform deployments do not require a \code{checkpointer=} argument in
\code{compile()}. The Platform provisions \code{AsyncPostgresSaver}
automatically and injects it into the compiled graph at runtime. Passing
\code{checkpointer=} to \code{compile()} in a Platform deployment is
silently ignored; the Platform's checkpointer always takes precedence.
Chapter~23~\S23.1 adds \code{EncryptedSerializer} wrapping a self-hosted
\code{AsyncPostgresSaver} for deployments where checkpoint data must not
be readable by the database administrator.
\end{notebox}

The full self-hosted stack---persistence, recovery, and dispatch---and its
managed equivalent are now both established. The remaining problem is
verification: no amount of unit testing with \code{InMemorySaver} can
confirm that the three layers compose correctly across real process
boundaries. That requires a test that kills actual processes and inspects
actual Postgres state between each termination.


%% ------------------------------------------------------------
\section{Correctness Is Proved by Killing the Process, Not Mocking It}
\label{sec:durable-full}
%% ------------------------------------------------------------

The design problem is the gap between unit-tested checkpointing behaviour
and actual process-boundary correctness. A unit test that mocks
\code{ainvoke} and asserts that \code{find\_resume\_checkpoint} returns
the expected tuple exercises the query logic in isolation but proves
nothing about the sequence that matters in production: that the checkpoint
is flushed to Postgres before the process exits, that a fresh process
reading the same \code{thread\_id} finds the correct row, and that the
graph resumes at the correct node rather than at \code{START}. The
governing principle is that the integration test must exercise an actual
process boundary---starting a subprocess, running it until it
self-terminates after writing a milestone, and starting a second subprocess
on the same \code{thread\_id}---with assertions on Postgres state between
each invocation.

The \code{DurableWorkflowState} schema carries a \code{stage} field
written by each node to name the current phase. This field plays a dual
role: it is the routing signal read by \code{route\_after\_recovery()}
when selecting the resume node, and it is the assertion target in the
test harness that confirms the correct milestone was committed before the
subprocess terminated. The three boolean flags---\code{research\_done},
\code{draft\_done}, \code{review\_done}---are written by their respective
nodes and provide a per-field view of execution progress independent of
the \code{stage} string, making assertion failures unambiguous: a failure
on \code{draft\_done} while \code{research\_done} is \code{True}
identifies exactly which phase was not committed.

\begin{lstlisting}[language=Python,
  caption={DurableWorkflowState and milestone node functions writing
           stage and boolean flags},
  label={lst:durable-state-nodes}]
# chapter22/state.py
from typing import TypedDict

class DurableWorkflowState(TypedDict):
    task:           str
    stage:          str   # "init"|"research"|"draft"|"review"|"done"
    research_done:  bool
    draft_done:     bool
    review_done:    bool
    final_output:   str

# chapter22/graph.py  (node functions)
from langchain_core.runnables import RunnableConfig
from chapter22.state import DurableWorkflowState

async def research_complete(
    s: DurableWorkflowState, _: RunnableConfig
) -> dict:
    return {"research_done": True, "stage": "research"}

async def draft_complete(
    s: DurableWorkflowState, _: RunnableConfig
) -> dict:
    return {"draft_done": True, "stage": "draft"}

async def review_complete(
    s: DurableWorkflowState, _: RunnableConfig
) -> dict:
    return {"review_done": True, "stage": "review"}

async def synthesise_node(
    s: DurableWorkflowState, _: RunnableConfig
) -> dict:
    return {"final_output": f"Done: {s['task']}", "stage": "done"}
\end{lstlisting}

Two decisions in Listing~\ref{lst:durable-state-nodes} deserve
explanation.
Node function names match the strings in \code{MILESTONE\_NODES}
exactly---\code{research\_complete}, \code{draft\_complete},
\code{review\_complete}---because \code{build\_durable\_graph()} passes
these function names to \code{b.add\_node()}: when Python registers
\code{b.add\_node("research\_complete",~research\_complete)}, the
\code{metadata["writes"]} key for that node's checkpoint will be the
string \code{"research\_complete"}, which matches the frozenset entry;
any deviation between the registered name and the frozenset entry
silently prevents recovery from ever finding a milestone.
\code{DurableWorkflowState} is a \code{TypedDict} rather than a Pydantic
\code{BaseModel}: checkpoint serialisation by \code{JsonPlusSerializer}
operates on plain dicts; a \code{TypedDict} round-trips through
serialisation without requiring \code{.model\_dump()} / \code{model\_validate()},
which is the correct type for values that live inside the Pregel
checkpoint store.

\begin{lstlisting}[language=Python,
  caption={build\_durable\_graph: RecoveryNode factory, conditional
           edge, and milestone node wiring},
  label={lst:durable-build-graph}]
# chapter22/graph.py  (continued)
from langgraph.graph import StateGraph, START, END
from langgraph.graph.state import CompiledStateGraph
from langgraph.checkpoint.postgres.aio import AsyncPostgresSaver
from chapter22.recovery import make_recovery_node, route_after_recovery

def build_durable_graph(
    checkpointer: AsyncPostgresSaver,
) -> CompiledStateGraph:
    b = StateGraph(DurableWorkflowState)
    b.add_node("recovery",          make_recovery_node(checkpointer))
    b.add_node("research_complete", research_complete)
    b.add_node("draft_complete",    draft_complete)
    b.add_node("review_complete",   review_complete)
    b.add_node("synthesise_node",   synthesise_node)

    b.add_edge(START, "recovery")
    b.add_conditional_edges(
        "recovery",
        route_after_recovery,
        {
            "research_complete": "research_complete",
            "draft_complete":    "draft_complete",
            "review_complete":   "review_complete",
            "synthesise_node":   "synthesise_node",
        },
    )
    b.add_edge("research_complete", "draft_complete")
    b.add_edge("draft_complete",    "review_complete")
    b.add_edge("review_complete",   "synthesise_node")
    b.add_edge("synthesise_node",   END)
    return b.compile(checkpointer=checkpointer)
\end{lstlisting}

Two decisions in Listing~\ref{lst:durable-build-graph} deserve
explanation.
The \code{add\_conditional\_edges} destination dict enumerates all four
possible return values of \code{route\_after\_recovery} explicitly:
LangGraph's static analysis validates at \code{compile()} time that every
string the routing function can return appears as a key in the dict; an
undeclared return value raises \code{ValueError} at compile time, not at
runtime, which surfaces routing bugs before the first invocation.
\code{make\_recovery\_node(checkpointer)} is called at graph-construction
time, closing over the \code{AsyncPostgresSaver} instance that was opened
in \code{lifespan\_checkpointer}; this is the same instance used by the
Pregel scheduler for its own checkpoint writes, ensuring that the
recovery scan and the graph's checkpoint writes share the same connection
pool and Postgres transaction semantics.

The subprocess worker installs an asyncio-native \code{SIGTERM} handler
before starting the graph run. The correct pattern for async processes
uses \code{loop.add\_signal\_handler()}, which schedules a coroutine on
the event loop when the OS delivers \code{SIGTERM}: because the callback
runs inside the loop, it can \code{await} the in-flight \code{aput} call
that the Pregel scheduler may have issued during the last node's checkpoint
flush. Without this protection, a \code{SIGTERM} delivered during
\code{aput} leaves a partial row in Postgres that subsequent reads may
interpret as a corrupt checkpoint.

\begin{warningbox}
Never call \code{asyncio.shield()} or schedule a coroutine directly
inside a synchronous signal handler registered with
\code{signal.signal()}. Synchronous signal handlers run in the main
thread between bytecodes---outside the event loop---and cannot schedule
or \code{await} coroutines. The result is a \code{RuntimeError: no
running event loop} that terminates the process before the checkpoint
flush completes. Use \code{loop.add\_signal\_handler()} to register an
async-compatible shutdown callback.
\end{warningbox}

\begin{lstlisting}[language=Python,
  caption={worker.py: asyncio-native SIGTERM handler shielding the
           in-flight checkpoint write; subprocess entry point},
  label={lst:durable-worker}]
# chapter22/worker.py
import argparse, asyncio, os, signal, sys

async def _shutdown(loop: asyncio.AbstractEventLoop) -> None:
    tasks = [t for t in asyncio.all_tasks(loop)
             if t is not asyncio.current_task()]
    for t in tasks:
        t.cancel()
    # Shield the gather from outer cancellation so in-flight
    # checkpoint writes (asyncpg aput) can complete cleanly.
    await asyncio.shield(
        asyncio.gather(*tasks, return_exceptions=True)
    )
    loop.stop()

async def main(thread_id: str, task: str) -> None:
    from chapter22.checkpointer import lifespan_checkpointer
    from chapter22.graph import build_durable_graph
    from chapter22.state import DurableWorkflowState
    from chapter22.settings import settings

    loop = asyncio.get_running_loop()
    loop.add_signal_handler(
        signal.SIGTERM,
        lambda: loop.create_task(_shutdown(loop)),
    )
    async with lifespan_checkpointer(settings.POSTGRES_DSN) as cp:
        graph = build_durable_graph(cp)
        cfg   = {"configurable": {"thread_id": thread_id}}
        await graph.ainvoke(
            DurableWorkflowState(task=task, stage="init",
                research_done=False, draft_done=False,
                review_done=False, final_output=""),
            config=cfg,
        )

if __name__ == "__main__":
    p = argparse.ArgumentParser()
    p.add_argument("--thread-id", required=True)
    p.add_argument("--task", default="Durability test")
    args = p.parse_args()
    asyncio.run(main(args.thread_id, args.task))
\end{lstlisting}

Three decisions in Listing~\ref{lst:durable-worker} deserve explanation.
\code{loop.create\_task(\_shutdown(loop))} is used inside the signal
handler lambda rather than \code{asyncio.ensure\_future()}: in Python~3.10+,
\code{asyncio.ensure\_future()} is deprecated in favour of
\code{loop.create\_task()}, and inside a callback registered with
\code{add\_signal\_handler} the loop is guaranteed to be running, making
\code{create\_task} the correct and unambiguous call. \code{asyncio.shield()}
wraps the \code{asyncio.gather()} inside \code{\_shutdown}: after task
cancellation is requested, the event loop may itself receive a second
cancellation from an outer scope (e.g., a timeout imposed by the OS or a
parent process sending \code{SIGKILL} after \code{SIGTERM}); without
\code{shield()}, that outer cancellation would abort the \code{gather}
before all tasks have flushed their final writes. The \code{thread\_id}
is passed as a CLI argument rather than an environment variable: the test
harness spawns multiple subprocesses with different \code{thread\_id} values
in the same test run; environment variables inherited across subprocesses
would require explicit clearing between spawns, whereas CLI arguments
are per-invocation and have no inheritance.

\begin{lstlisting}[language=Python,
  caption={Three-subprocess test harness: per-phase milestone assertion
           via alist() and final completion assertion},
  label={lst:durable-test}]
# chapter22/test_three_restarts.py
import asyncio, subprocess, sys, pytest
from langgraph.checkpoint.postgres.aio import AsyncPostgresSaver
from chapter22.recovery import MILESTONE_NODES
from chapter22.settings import settings

THREAD_ID = "durability-test-001"

async def assert_milestone(dsn: str, thread_id: str,
                            expected: str) -> None:
    async with AsyncPostgresSaver.from_conn_string(dsn) as cp:
        await cp.setup()
        async for tup in cp.alist(
            {"configurable": {"thread_id": thread_id}}
        ):
            writes = tup.metadata.get("writes") or {}
            if MILESTONE_NODES.intersection(writes.keys()):
                node = next(k for k in writes if k in MILESTONE_NODES)
                assert node == expected, \
                    f"Expected {expected!r}, got {node!r}"
                return
    pytest.fail(f"No milestone checkpoint found for {thread_id!r}")

def _phase(stop_after_signal: bool = False) -> subprocess.CompletedProcess:
    cmd = [sys.executable, "chapter22/worker.py",
           "--thread-id", THREAD_ID, "--task", "Durability test"]
    proc = subprocess.run(cmd, timeout=120)
    # 0 = clean exit; -15 = SIGTERM before handler registered.
    assert proc.returncode in {0, -15}, \
        f"Unexpected exit code: {proc.returncode}"
    return proc

@pytest.mark.integration
def test_three_restarts() -> None:
    milestones = [
        "research_complete",
        "draft_complete",
        "review_complete",
    ]
    for expected in milestones:
        _phase()
        asyncio.run(assert_milestone(
            settings.POSTGRES_DSN, THREAD_ID, expected
        ))
    # Fourth run: no interruption; assert full completion.
    proc = subprocess.run(
        [sys.executable, "chapter22/worker.py",
         "--thread-id", THREAD_ID, "--task", "Durability test"],
        timeout=120,
    )
    assert proc.returncode == 0
    async def _final():
        async with AsyncPostgresSaver.from_conn_string(
            settings.POSTGRES_DSN
        ) as cp:
            await cp.setup()
            snap = await cp.aget(
                {"configurable": {"thread_id": THREAD_ID}}
            )
            vals = snap.checkpoint["channel_values"]
            assert vals["final_output"] != "", \
                "final_output must be written after full completion"
            assert vals["stage"] == "done"
    asyncio.run(_final())
\end{lstlisting}

Three decisions in Listing~\ref{lst:durable-test} deserve explanation.
\code{@pytest.mark.integration} marks the test for selective execution:
this test requires a live Postgres and Redis instance and sends actual
\code{SIGTERM} signals; running it in a standard unit-test suite without
infrastructure would fail or block; the \code{integration} marker allows
\code{pytest -m "not integration"} to exclude it from CI pipelines that
run unit tests only. The final assertion uses
\code{cp.aget(\{"configurable":~\{"thread\_id":~...\}\})} rather than
\code{alist()}: \code{aget()} returns the single latest
\code{CheckpointTuple} without iteration, which is the correct call when
the caller needs only the most recent state and not the full history;
using \code{alist()} here and taking the first result would be
equivalent but semantically misleading. \code{assert\_milestone} calls
\code{cp.setup()} inside the assertion helper because the helper opens
its own \code{AsyncPostgresSaver} context independent of any running
worker process; calling \code{setup()} here is safe---it is idempotent---
and removes the requirement that the schema already exists before the
assertion runs, which could cause the first assertion call to fail with
a \code{UndefinedTableError} in a freshly provisioned test database.

The durability stack proved here---\code{AsyncPostgresSaver} checkpointing
to Postgres, \code{RecoveryNode} resuming from the highest milestone,
Celery deferring invocations through a durable task queue---provides the
infrastructure foundation that Chapter~23 builds upon when adding
encrypted checkpoint serialisation and multi-tenant thread isolation,
where the correct
